\section{Open Problems}
\label{sec:open_problems}


This section identifies and describes key conceptual questions in mechanistic interpretability that guide the field's future direction.






\subsection{When Is Mechanistic Understanding Necessary?}

Section~\ref{subsec: mech interp necessary} discusses cases in which mechanistic understanding is necessary in principle, such as when the goal is to characterise the internal computations that produce a model's behaviour. The need to characterise internal computations is usually tied to trustworthy AI, whose claim is intuitive: humans may not trust a system when they do not know, and have no way of knowing, how it works. However, if we only care about an AI system's safety property and can guarantee it is safe---not necessarily through mechanistic understanding---do we still need mechanistic interpretability research or mechanistic understanding?

For example, sleeper agents' deceptive behaviours hidden from standard behavioural evaluations~\citeem{hubinger2024sleeperagentstrainingdeceptive} can leave detectable signatures in the model's internal representations, and \citetem{macdiarmid2024sleeperagentprobes} showed that this behaviour can be detected from model activations. This demonstrates the utility of looking into the internals. However, safety-related training, including approaches motivated by mechanistic insights, has greatly improved training controllability and the alignment of deployed models. Adversaries may now find it harder to implement such a backdoor. Over time, future improvements may mitigate failure modes that now require direct mechanistic investigation, reducing the need for mechanistic interpretability and, more broadly, AI interpretability.

The AI safety community has not yet agreed on whether this will happen. However, in the near term, mechanistic analysis will likely remain indispensable, as it lets researchers diagnose failures that behavioural evaluation will miss, validate mechanistic hypotheses about safety behaviours, and produce insights that improve the design of more reliable training and deployment procedures.

\subsection{Challenges in Evaluation and Benchmarking}
Because widely accepted standards for benchmarks and evaluation are lacking, evaluating mechanistic interpretability methods remains hard. Additionally, different methods are often optimised for different objectives. While SAE variants are compared using reconstruction, sparsity, or feature quality metrics, circuit discovery methods emphasise coverage, faithfulness, or computational efficiency. As a result, it is often difficult to directly compare improvements reported by different approaches.

Section~\ref{sec:evaluation-framework} proposes a set of evaluation dimensions, but turning them into shared benchmarks remains open. Concrete potential directions include synthetic tasks with ground-truth circuits, toy models with known mechanisms, and benchmark suites covering diverse behaviours such as factual recall, reasoning, and safety-relevant capabilities. Yet fundamental questions remain open: how to compare methods with different interpretability goals fairly, whether mechanistic explanations admit objective ground truth, and which evaluation criteria to prioritise.

\subsection{Incompleteness, Scalability, Generalisability, and Adversarial Robustness}
\label{subsec: scalability and adversarial robustness}






Mechanistic methods, the insights they uncover, and their downstream use face several critical challenges. Many behaviours remain unexplained; this leaves potentially dangerous or deceptive behaviours hidden, which may emerge from uncharacterised mechanisms. 
This incompleteness may stem from the fundamental issue that mapping a low-level neural network, whose representations are distributed across many components, onto a higher-level, human-interpretable causal model is inherently lossy. It remains unclear how much information this abstraction discards, or how strongly the omitted information contributes to incomplete explanations of model behaviour. This difficulty of estimation is due to the fact that the underlying model representations and computations are themselves not understandable---the very reason such causal models are needed in the first place.

Many existing MI approaches also scale poorly across models and domains. Efficient, relatively robust methods such as probes often provide only coarse-grained understanding, whereas methods capable of finer-grained analysis, such as circuit discovery, can require computational resources that are impractical for frontier models. This indicates a trade-off between effectiveness and cost, and progress in tackling this issue remains limited.

Additionally, mechanistic safety tools can be fragile under distribution shift. For instance, monitoring methods or steering vectors may perform reliably in distribution but degrade substantially when the input distribution shifts~\citeest{tan2024analysing}, raising concerns about their robustness and generalisability in real-world deployment.

Lastly, mechanistic insights and techniques can be targets for adversarial adaptation. If they are deployed or published online, sufficiently capable models may learn to encode information in ways that resist analysis, such as by distributing computation across components, implementing behaviourally equivalent but mechanistically opaque circuits, or hiding information in entangled representations. This possibility is analogous to evaluation awareness~\citeem{anthropic2025claude-sonnet-4-5} and raises the prospect of an interpretability arms race with no clear endpoint, in which stronger interpretability methods create incentives for models to become correspondingly harder to inspect.

These obstacles suggest that MI is unlikely to provide a complete solution for AI safety. Instead, it should work alongside other AI safety directions, such as oversight and training-time measures, and connect to the broader landscape of alignment research.

\subsection{Mechanistic Understanding and Alignment}
\label{sec:mechanistic-alignment}

MI provides tools for behavioural control. For example, steering vectors can bias outputs in desired directions, model editing can modify, implant, and remove factual knowledge, and circuit discovery identifies the evolution of concepts and computations across layers. 
Yet these approaches do not directly tackle alignment specification. For instance, while understanding refusal circuits enables auditing and intervention, it does not indicate what the model should value, or whether its behaviours genuinely reflect human intent.

Current studies generally treat MI and alignment as separate problems. Mechanistic methods support verification and control, while other approaches, such as constitutional AI~\citeem{bai2022constitutional} and RLHF~\citeem{bai2022training}, focus on specifying desired behaviours. Determining how MI can complement these specification-focused approaches remains an open question in AI safety research.

\subsection{Interpretation vs. Truth: The Ontology Question}

\label{sec:ontology}



A foundational MI question, which is both technical and conceptual, is: what does it mean to mechanistically understand a neural network? Marking a feature as a ``borzoi detector'', for instance, or labelling an attention head component as ``a previous-token copier'' mirrors possibly subjective human choices. This means multiple mechanisms can explain one behaviour, each consistent with observed causal relationships.

MI research often assumes how human-interpretable descriptions should map onto causal mechanisms. But these assumptions face several obstacles. As mentioned before, one of these is subjectivity. This includes choosing the level of abstraction; for instance, a circuit could be explained at the level of attention routing or low-rank decomposition, and the choice depends on subjective priorities. As such, while researchers may obtain explanations that are intuitive and fit external observations, they may not fully represent the actual internal computations~\citeest{hewitt2025position}.

Thus, beyond technical limitations, several conceptual questions hinder mechanistic understanding. Other questions include: How can we assess confidence in a mechanistic explanation if multiple descriptions can explain the same behaviour? Are multiple valid interpretations acceptable, or should researchers expect a unique mechanistic description? Resolving these questions will help researchers clarify the limits and goals of MI research.

\subsection{The Proxy Fidelity Problem}
\label{sec:proxy-fidelity}















Several recent interpretability approaches construct more interpretable proxy models that approximate the target model's internal computations. Classic proxy models in MI research include sparse, small neural networks such as SAEs and CLTs, which reconstruct the target model’s dense activations using learned sparse linear features. However, while these approaches have greatly enhanced mechanistic analysis, they introduce a new obstacle: how can we ascertain whether the proxy truly learns the same mechanisms as the original model, rather than shallowly reproducing behaviour? This concern is realistic because two neural networks can produce the same output while using different internal algorithms. That is, a proxy model may find features and circuits that do not exist in the original model.

Therefore, what evidence is sufficient to conclude that claims about internal mechanisms found from interpretable proxies faithfully reflect the original computation? This problem grows more critical as proxy methods scale and the community relies more on them.

\paragraph{Why This Matters for Safety.} For example, when researchers find no ``deception circuits'' in the proxy models, can they claim that the target model has therefore ``passed the test'' and is safe? Making this claim requires that these mechanisms do not exist in the target model either. However, as discussed above, the absence of a mechanism in the proxy model does not imply its absence in the target model, since they may compute the input differently while producing the same output. The conclusion we can draw is only something like ``we did not find the deception mechanism in the proxy model'', which is far from a useful guarantee in safety-critical scenarios.

\paragraph{Potential Approaches.}

Several directions can validate and improve proxy fidelity. One promising direction is interchange intervention testing and training~\citeest{das, pmlr-v162-geiger22a}. This technique was developed to compare and align a neural network with a high-level causal abstraction. The proxy models and their circuits can be viewed as a causal abstraction of the target model, where the nodes (human-interpretable variables) are the proxy models' features. This framing not only lets us test alignment between the mechanisms identified in proxy models and the target model, but also provides an optimisation objective for training faithful proxies. Additionally, a theoretical framework that upper-bounds how much internal computation can differ between the target model and proxy model under a given reconstruction performance can also be a promising direction.

\subsection{Transferability Beyond Transformers}
\label{sec:non-transformer}

The methods surveyed in this paper have been developed and evaluated predominantly in transformer-based autoregressive language models. Although the six research orientations discussed in Section~\ref{sec:orientations} are largely architecture-agnostic in their aims, their classic methods and representative works focus almost entirely on transformers. Even though foundational circuits work~\citeem{olah2020zoom} originated in convolutional networks (InceptionV1; \citeestalp{szegedy2015going}), demonstrating interpretable features of visual neurons, many of the techniques that define modern mechanistic interpretability, which include feature extraction and explanation, model editing, and circuit discovery, were developed specifically for transformer residual streams, MLPs, and attention mechanisms. Thus, an important question arises: are mechanistic interpretability concepts, tools, and findings specific to transformers, or do they reflect more general properties of neural computation across model architectures? Below, we discuss several representative models that are substantially different from transformer-based autoregressive models.

\begin{itemize}
    \item \textbf{Convolutional neural networks (CNNs):} CNNs are a representative vision model architecture~(\citeemalp{oshea2015introductionconvolutionalneuralnetworks}; \citeestalp{NIPS2012_c399862d}). Their fundamental building blocks are convolutional and pooling layers rather than MLP blocks and self-attention. Transformer-specific circuit formulations thus do not transfer directly to CNNs, although circuit analysis itself originated there~\citeem{olah2020zoom}. Convolutions' spatial locality may yield more naturally modular circuits than transformers under suitable formulations, but large-scale studies of superposition, feature extraction, and circuit discovery in modern CNNs and their variants remain limited.

    \item \textbf{State-space and recurrent models:} Architectures such as Mamba~\citeest{gu2024mamba} and RWKV~\citeest{peng-etal-2023-rwkv} replace standard self-attention with recurrent computations that update a hidden state over time. They can be seen as modern recurrent neural network architectures~\citeest{elman1990finding}. Specifically, Mamba uses input-dependent state transitions to selectively retain or forget information, while RWKV uses a recurrence based on linear attention. This new paradigm raises interpretability questions that previous transformer interpretability studies have not addressed, such as how information is stored and updated in recurrent states, whether these representations exhibit certain properties such as superposition, and how causal effects can be traced across state updates.

    \item \textbf{Diffusion models:} Diffusion models~\citeest{pmlr-v37-sohl-dickstein15} apply iterative denoising processes to generate images or tokens. They may still use transformers as backbones, but their computational structure differs fundamentally from autoregressive generation. For example, many current MI studies use the predicted next token as a sign of model behaviour (e.g., ``yes'' or ``no'' when the input is a malicious query), while diffusion-based language models generate multiple tokens in different positions at a time to iteratively complete the response. This difference means many interpretability techniques and insights do not directly apply. While recent work has begun investigating attention patterns and internal representations in diffusion models~(\citeestalp{tinaz2025emergence}; \citeemalp{wang2026dlmscopemechanisticinterpretabilitydiffusion}), mechanistic understanding in this paradigm remains lacking.
\end{itemize}

Beyond these major architectural changes, researchers have also proposed variants of conventional transformer components. For example, mixture-of-experts (MoE) layers~\citeest{lepikhin2021gshard, dai-etal-2024-deepseekmoe} replace standard MLP blocks with multiple smaller MLPs, called ``experts''. For each input activation, a router selects a subset of experts based on the activation and combines their outputs using learned routing weights. Many modern transformer-based language models have adopted this architecture~\citeem{jiang2024mixtralexperts, deepseekai2025deepseekv3technicalreport, yang2025qwen3technicalreport, openai2025gptoss120bgptoss20bmodel, kimiteam2026kimik2openagentic}.

However, MoE architectures introduce technical challenges for mechanistic analysis. For example, activation patching at an intermediate MoE layer may change which experts are selected in downstream MoE layers. This makes it difficult to distinguish the patched activation’s effects on model output from those mediated by changes in downstream expert routing. Furthermore, MoE layers can contain many experts, ranging from 8~\citeem{jiang2024mixtralexperts} to 256~\citeem{deepseekai2025deepseekv3technicalreport} or more~\citeem{deepseekai2026deepseekv4highlyefficientmilliontoken}. Although only a subset of experts is active for each input, the large number of possible expert pathways increases the search space for circuit discovery, making circuit discovery harder to use in modern MoE-based models.

Simultaneously, some recent work suggests that at least a few mechanistic interpretability techniques can transfer outside of transformers. Steering vectors, tuned lenses, knowledge editing, and circuit analysis evaluations have all been adapted, with varying success, to state-space models such as RWKV and Mamba, as well as to vision-language models~(\citeestalp{paulo2024does, sharma2024locating, haon2025mechanistic, neo2025towards}; \citeemalp{ensign2024investigating}). Researchers have also conducted theoretical analyses aiming to clarify the computational relationships between transformers and other architectures. Some initial results suggest that several findings in transformers may reflect more general, transferable properties of neural networks~\citeest{zimerman2025explaining, dao202mamba2, peng2025rwkv, ali-etal-2025-hidden}. Sorting out which aspects of mechanistic interpretability depend on architecture and which reflect broader principles of neural computation thus appears to be a promising research direction. On the other hand, whether current research orientations and associated techniques remain relevant beyond transformers will also show how broadly mechanistic interpretability research can apply to future systems.
